
% ===================================================================
\section{Ensemble models}\label{app:ensemble}
% ===================================================================
The ensemble maps the donor pool to a synthetic Brent series through five estimators, each
with a different inductive bias. Where the data underdetermine the counterfactual the five
disagree in structured ways rather than failing in the same direction, and we read that spread
as model uncertainty. The convex and elastic-net members are sparse and stay inside or close to
the donor envelope, the augmented member adds a ridge layer that permits limited extrapolation,
gradient boosting captures nonlinear interactions but cannot extrapolate past any donor's
training range, and the Bayesian ridge retains every donor under smooth shrinkage. We fit each
member on its own, combine them by the cross-sectional median, and report the per-model gaps in
the body. This appendix gives the estimator definitions, the hyperparameter search and selected
settings, the robustness grid that varies the main modelling choices, and a note on the Bayesian
member and its relation to the full BSTS model.

% ===================================================================
\subsection{The five estimators}\label{app:formulas}
% ===================================================================

\paragraph{Convex SCM \citep{abadie2010}.} A simplex-constrained weighted average of donors,
\begin{equation}
\hat y_t = \sum_{j=1}^{J} w_j\, y_{j,t}, \qquad w_j \ge 0, \quad \sum_j w_j = 1 .
\label{eq:convex}
\end{equation}
The weights read like portfolio shares and the synthetic cannot leave the donor envelope. The
simplex constraint behaves like an $L_1$ ball, so the optimiser drives most $w_j$ to exactly
zero: the fit is mechanically sparse and leans on the few donors whose convex combination best
reproduces pre-period Brent.

\paragraph{Augmented SCM \citep{benmichael2021}.} The convex base is kept and a ridge layer is
added that extrapolates each donor's contribution beyond the simplex,
\begin{equation}
\hat y_t = \underbrace{\sum_j w_j\, y_{j,t}}_{\text{convex base}}
\;+\; \underbrace{\sum_j \eta_j\, y_{j,t}}_{\text{ridge correction}},
\qquad \eta = (X^\top X + \lambda I)^{-1} X^\top r,
\label{eq:ascm}
\end{equation}
where $r$ is the pre-period residual left by the convex base and $X$ the donor matrix. The
effective weight is $w_j+\eta_j$; the $\eta_j$ are unconstrained, which is what permits limited
extrapolation outside the donor hull. Large $\lambda$ shrinks the ridge layer toward zero and
recovers convex SCM; small $\lambda$ gives it more freedom.

\paragraph{Elastic net \citep{doudchenko2016,zouhastie2005}.} A signed linear regression with a
combined $L_1+L_2$ penalty,
\begin{equation}
\hat y_t = \beta_0 + \sum_j \beta_j\, y_{j,t}, \qquad
\min_\beta \;\tfrac{1}{2N}\lVert y - X\beta\rVert_2^2
\;+\; \alpha\Big[\rho\lVert\beta\rVert_1 + \tfrac{1-\rho}{2}\lVert\beta\rVert_2^2\Big].
\label{eq:enet}
\end{equation}
The $L_1$ term drives some coefficients to exactly zero (selection); the $L_2$ term shrinks the
rest and stabilises the fit under donor collinearity. Here $\alpha$ is the overall penalty
strength and $\rho$ the $L_1$ share ($\rho=0$ pure ridge, $\rho=1$ pure lasso).

\paragraph{Gradient boosting (XGBoost) \citep{chenguestrin2016}.} A sum of $K$ shallow
regression trees fit sequentially to residuals,
\begin{equation}
\hat y_t = \sum_{k=1}^{K} f_k(\mathbf{y}_{\cdot,t}),
\qquad f_k \in \{\text{trees of depth} \le d\},
\label{eq:xgb}
\end{equation}
with $L_2$ regularisation on leaf weights and a learning rate $\eta$. It is not a closed-form
regression: at each split the builder picks the donor and threshold that most reduce the
residual. It captures nonlinear interactions but, by construction, cannot extrapolate outside
the training range of any donor --- the one member that brackets the extrapolating linear
models from below.

\paragraph{Bayesian ridge \citep{brodersen2015}.} The OLS regression equation with a Gaussian
prior on the coefficients and Gamma priors on the precisions,
\begin{equation}
y_t = \beta_0 + \sum_j \beta_j\, y_{j,t} + \varepsilon_t,\quad
\varepsilon_t\sim\mathcal{N}(0,1/\alpha);\qquad
\beta_j\mid\lambda \sim \mathcal{N}(0,1/\lambda),\quad \lambda,\alpha\sim\text{Gamma}(\cdot).
\label{eq:bridge}
\end{equation}
The posterior over $\beta$ is conjugate and closed-form, so $\hat y_t$ carries a native
posterior standard deviation. The Gaussian prior shrinks coefficients smoothly toward zero but
never sets any to exactly zero, so --- unlike the convex and elastic-net members --- it performs
no selection and retains every donor. This implements only the regression-on-donors term of the
BSTS model of \citet{brodersen2015}; the trend, seasonal, and spike-and-slab components it omits,
and why the proxy is adequate at this sample size, are discussed in Appendix~\ref{app:bsts}.

% ===================================================================
\subsection{Hyperparameter grids and selected configurations}\label{app:grids}
% ===================================================================
Hyperparameters are chosen by the pooled validation RMSE of the expanding-window walk-forward
cross-validation described in the body (each event tuned on its own pre-window), with the final
model refit on the full pre-window. Table~\ref{tab:grids} lists the search grids as run in the
fitting pipeline, restricted to the one or two knobs that materially change each model's fit; the
selection (winner's-curse) bias of \citet{cawley2010} grows with the number of configurations
tested, which is why the grids stay narrow. Convex SCM has no tuned hyperparameter: it draws 120
importance matrices $V$ from a Dirichlet, re-fits weights per $V$, and retains the $V$ with lowest
pre-period RMSPE.

\begin{table}[h]
\centering
\caption{Hyperparameter search grids as run in the fitting pipeline (the grids that produced the
reported fits). Selection is by pooled walk-forward validation RMSE, per event; Convex SCM is
untuned. XGBoost uses early stopping (20 rounds) on each fold's validation window, with the
median best iteration carried to the final refit.}
\label{tab:grids}
\small
\begin{tabular}{@{}lp{10.8cm}@{}}
\toprule
Estimator & Search grid \\
\midrule
Convex SCM     & None tuned; 120 Dirichlet draws of the importance matrix $V$, best by pre-period RMSPE.\\
Augmented SCM  & Ridge penalty $\lambda \in \{0.01,\,0.1,\,1,\,10,\,100\}$ (5).\\
Elastic net    & $\alpha \in \{0.001,\,0.01,\,0.05,\,0.1,\,0.5\}$ $\times$ $\rho \in \{0.2,\,0.5,\,0.8\}$ (15).\\
XGBoost        & \texttt{max\_depth} $\in \{2,3,4\}$ $\times$ \texttt{learning\_rate} $\in \{0.001,\,0.005,\,0.01,\,0.03,\,0.1\}$ $\times$ \texttt{reg\_lambda} $\in \{1,10\}$ (30); \texttt{n\_estimators}\,$=$\,1000 with early stopping after 20 rounds, \texttt{gamma}\,$=$\,0.1.\\
Bayesian ridge & Weight precision $\lambda \in \{10^{-8},10^{-6},10^{-4},10^{-2},1\}$ $\times$ noise precision $\alpha \in \{10^{-8},10^{-6},10^{-4}\}$ (15).\\
\bottomrule
\end{tabular}
\end{table}

\emph{Selected configurations} (lowest pooled validation RMSE, read from the saved fits).
\textbf{Russia:} ASCM $\lambda=0.1$; elastic net $\alpha=0.001,\;\rho=0.2$; XGBoost depth~2,
learning rate~$0.1$, $\lambda_{\text{reg}}=1$ (early-stopped near 32 trees); Bayesian ridge
$\lambda=1,\;\alpha=10^{-4}$. \textbf{Hormuz:} ASCM $\lambda=10$; elastic net
$\alpha=0.01,\;\rho=0.2$; XGBoost depth~3, learning rate~$0.1$, $\lambda_{\text{reg}}=10$
(early-stopped near 10 trees); Bayesian ridge $\lambda=1,\;\alpha=10^{-8}$. The two events select
genuinely different regularisation strengths, as expected: Russia's larger convex residual asks
the ASCM ridge layer for more freedom (small $\lambda$), while Hormuz's tighter convex fit
tolerates heavy shrinkage. These are the settings behind the per-model gaps reported in the body.

% ===================================================================
\subsection{Robustness grid}\label{app:grid}
% ===================================================================
The headline uses the preferred specification (shared 19-donor pool, preferred pre-window,
ensemble median). The full grid varies (i) the pre-window (preferred, extended, and narrow
specifications); (ii) the donor pool (shared 19, permissive 26 for Russia, full 32); and
(iii) the model (each of the five plus the median). For Russia the donor-pool sweep is
itself a diagnostic: the gap is expected to \emph{shrink} monotonically from the clean
19-donor pool to the contaminated 32-donor pool, because Russia-treated donors get inflated
by the same shock and pull the synthetic up. If this monotonic pattern holds, the
SUTVA-driven exclusion logic is empirically supported.

% ===================================================================
\subsection{Bayesian ridge and the full BSTS model}\label{app:bsts}
% ===================================================================
The fifth ensemble member implements only the regression-on-donors term $\beta'x_t$ of the
Bayesian structural time-series model of \citet{brodersen2015}, with a conjugate
Gaussian--Gamma prior. It omits the local-linear \emph{trend}, the \emph{seasonal}
component, the \emph{spike-and-slab} donor selection, and the autocorrelation-aware
posterior over the full state trajectory. At our sample size, with no strong weekly
seasonality in daily oil prices, the regression term carries most of the signal, so the
proxy is an adequate Bayesian-flavoured diversifier; but its credible intervals are Gaussian
and ignore autocorrelation, and its donor ``importances'' are standardised-coefficient
magnitudes, not true posterior inclusion probabilities. We therefore call it a Bayesian
ridge regression rather than a BSTS, and read its placebo $p$-values as rank statistics
rather than formal probabilities (Table~\ref{tab:placebo}). A full BSTS implementation is the
natural route to autocorrelation-aware credible intervals in a future iteration.
